\chapter{La demostración}

\inicio{A}{} las once en punto, el Consejo del Archivo volvió a reunirse en la sala de las paredes de cristal fundido. Esta vez había más gente. Se había corrido la voz de que la doctora Calderón iba a desafiar al Gremio de la Llave, y en el Archivo de Ceres no ocurría nada tan emocionante desde la avería del ascensor de 2179.

La caja de llaves seguía en el centro de la mesa. Inés no la miró. Pidió una pizarra, la consiguió tras alguna resistencia, y en veinte minutos explicó el talonario, el gesto único de arrancar, el procedimiento de cambiar un cartel y la recarga con salida temprana. Siete la ayudó con los ejemplos. Beto Arce, desde un rincón, asentía con tanto entusiasmo que en dos ocasiones la Directora le pidió que parara.

Cuando terminó, Helga Sørensen guardó silencio un momento.

---Es ingenioso ---dijo por fin---\nobreak. Lo reconozco. Pero el Gremio no adopta procedimientos ingeniosos, doctora; adopta procedimientos seguros. Toda su construcción descansa en una suposición: que el talonario nunca baja de cero. Si baja una sola vez, el mecanismo da la vuelta, el cartel se llena de tickets fantasma y estamos peor que antes, porque al menos antes el fallo era raro. Usted me ha enseñado ejemplos con cuatro lectores. El Archivo tiene sesenta y cinco mil. ¿Puede demostrarme que, con sesenta y cinco mil lectores intercalándose de todas las formas posibles, el talonario nunca baja de cero?

---Sí ---dijo Inés.

---¿Sin ejemplos?

---Sin ejemplos. Los ejemplos convencen; las demostraciones obligan. ---Borró la pizarra---\nobreak. Necesito dos definiciones y dos afirmaciones. Las definiciones son gratis. Las afirmaciones se las demuestro.

\escena
---Primera definición ---dijo Inés, y la escribió---\nobreak. Un lector comete una \emph{ofensa} si arranca un ticket cuando el talonario tiene $P$ tickets o menos. Es decir: si mete la mano en la zona roja. Quien comete una ofensa queda, por el procedimiento, obligado a recargar antes de marcharse.

---Una palabra fuerte para algo tan pequeño ---comentó la Directora.

---Es una palabra de abogado, y la demostración es un juicio. Segunda definición: llamemos $O$ a la lista de ofensores desde la última vez que el talonario estuvo por encima de la zona roja. Cada vez que alguien ofende, se añade a $O$. Cada vez que alguien pone el talonario por encima de $P$ ---sea cambiando el cartel, sea recargando con éxito---\nobreak, $O$ se vacía.

---¿Una lista con repetidos? ---preguntó Loew, que esa mañana estaba sorprendentemente despierto---\nobreak. ¿Puede aparecer el mismo lector dos veces?

---En principio sí, y por eso los matemáticos la llamarían un \emph{multiconjunto}: una bolsa en la que caben bolas repetidas. Pero precisamente lo primero que voy a demostrar es que eso no ocurre.

Escribió en la pizarra:

\begin{pizarra}
\textsc{Afirmación 1.} Ningún lector está en $O$ más de una vez.

\smallskip
\textsc{Afirmación 2.} En todo momento, los tickets del talonario son al menos $P - |O|$, donde $|O|$ es el número de elementos de $O$.
\end{pizarra}

---Si las dos son ciertas ---dijo---\nobreak, el resto es aritmética de escuela. Hay como mucho $P$ lectores. Por la primera, cada uno está en $O$ como mucho una vez, así que $|O|$ no pasa de $P$. Por la segunda, el talonario es al menos $P - |O|$, que es al menos $P - P$, que es cero. El talonario nunca es negativo.

Hubo un silencio en la sala. Mendieta, que no había entendido nada desde la palabra «multiconjunto», tuvo sin embargo la impresión de que algo importante acababa de pasar.

---Demuestre la primera ---dijo la Directora.

\escena
---Un lector entra en $O$ por una sola vía: ofendiendo, es decir, arrancando un ticket en la zona roja. Y nadie más que él puede meterlo en $O$: es su propia acción. ---Inés trazó una flecha---\nobreak. Cuando ofende, su lectura no termina: tiene que recargar. Así que, antes de poder ofender otra vez, tiene que terminar su recarga. La pregunta es: ¿puede terminar su recarga sin que $O$ se haya vaciado?

---Veamos ---siguió---\nobreak. Hay tres formas de salir de la recarga. Una: su cotejo funcionó. Entonces él mismo puso el talonario por encima de $P$, por R-1, y $O$ se vació. Dos: vio otro cristal en el cartel. Entonces alguien cambió el cartel, el nuevo talonario tiene $C$ tickets, más que $P$, y $O$ se vació. Tres: vio más tickets de los que esperaba. Entonces, por R-2, el talonario saltó por encima de $P$ en algún momento, y $O$ se vació.

---En los tres casos ---dijo Siete, desde su sitio---\nobreak, entre dos ofensas del mismo lector siempre hay un vaciado de $O$.

---Por tanto, nadie está dos veces. Primera afirmación demostrada.

La Directora no dijo nada. Inés lo tomó como un permiso para seguir.

---Segunda. Miremos lo que ocurre desde el último vaciado de $O$. Justo después del vaciado, el talonario está por encima de $P$ y $O$ está vacío: los tickets son más que $P - 0$. Se cumple. Desde entonces, ¿cómo puede haber cambiado el talonario? Por R-2, solo bajando de uno en uno; porque si hubiera subido, habría saltado por encima de $P$, y $O$ se habría vaciado otra vez, y estaríamos mirando desde un vaciado posterior.

Dibujó una escalera descendente.

---Cada bajada es un ticket arrancado. Las bajadas que ocurren por encima de $P$ dejan el talonario en $P$ o más: no estropean nada. Y cada bajada que ocurre desde $P$ o menos es, por definición, una ofensa: baja el talonario en uno y suma uno a $|O|$. Así que el talonario está, como mucho, $|O|$ escalones por debajo de $P$. Es decir: los tickets son al menos $P - |O|$. Segunda afirmación demostrada.

Dejó la tiza sobre la mesa.

---Y con las dos, el talonario nunca es negativo. Con cuatro lectores o con sesenta y cinco mil, se intercalen como se intercalen.

\escena
Fue Loew quien rompió el silencio, con una risa suave.

---Ahora entiendo por qué la zona roja empieza en $P$ y no en cero ---dijo---\nobreak. Es el colchón exacto. En el peor caso, todos los lectores del Archivo ofenden a la vez, cada uno una sola vez, antes de que ninguna recarga termine. Y el colchón de $P$ tickets aguanta exactamente eso. Ni uno más, ni uno menos.

---Exacto, profesor. En los ensayos de esta mañana, con cuatro lectores, el talonario llegó a cero una vez: los cuatro habían ofendido a la vez. Tocó la cota sin cruzarla. La demostración no solo es correcta; es ajustada.

---Doctora ---dijo Siete, y había en su voz un temblor que Inés no le conocía---\nobreak, me he dado cuenta de algo mientras usted hablaba. Todos los lectores que están en $O$ están, en ese momento, a mitad de una recarga, porque quien ofende no se va sin recargar. Y por la segunda afirmación, el talonario solo puede estar por debajo de $P$ si $O$ no está vacío. Por tanto\ldots

---Dilo.

---Por tanto, si el talonario está en zona roja, \emph{siempre hay alguien recargándolo}. No puede existir un momento en que el talonario esté bajo y nadie se ocupe de él.

Inés lo miró un momento largo.

---Eso no lo había pensado yo ---dijo---\nobreak. Apúntatelo. Es tuyo.

Siete no contestó, pero Beto aseguró después que se le iluminaron los ojos un tono más de lo reglamentario.

\escena
---Muy bien ---dijo la Directora, que no había cambiado de expresión en todo el rato---\nobreak. El talonario no baja de cero. Ahora hábleme del coste. ¿Cuánto tarda un lector en leer un cartel?

---Casi siempre, un gesto ---dijo Inés---\nobreak. Arranca su ticket y se va. Solo si deja el talonario en zona roja tiene que recargar, y la recarga, como Siete demostró esta mañana, cuesta como mucho $P$ intentos de cotejo. Y no siempre es el mismo lector: cada vuelta de más en la recarga de alguien requiere que otro lector ofenda en ese mismo cartel, y cada lector solo puede ofender una vez. Si solo tres robots usan un cartel, la recarga cuesta como mucho unos tres intentos, aunque el Archivo tenga sesenta mil.

---¿Y con qué frecuencia hay que recargar?

---Con un talonario de diecinueve casillas, podemos cargar del orden de medio millón de tickets. Tras una recarga, el talonario queda casi lleno, y no vuelve a hacer falta recargar hasta bajar de sesenta y cinco mil. Del orden de cuatrocientas cincuenta y ocho mil lecturas entre recargas. Dos de cada millón de lecturas recargan. Las demás son un solo gesto.

---¿Y cambiar el cartel?

---Tres gestos: apuntar en el nuevo, trocar, restar en el viejo. Siempre tres, haya cuantos lectores haya. Sin esperas.

La Directora tamborileó con los dedos sobre la caja de llaves.

---Ha dicho usted «sesenta y cinco mil lectores como máximo» varias veces. ¿Qué pasa si hay más?

Inés no esquivó la pregunta.

---Entonces la demostración deja de valer. Si más de $P$ lectores usan a la vez el mismo cartel, $|O|$ podría superar a $P$, el talonario podría bajar de cero y tendríamos tickets fantasma. Es un requisito estricto, Directora, no una recomendación. El censo del Archivo tiene que garantizar que nunca hay más de $P$ lectores. Por suerte, el censo es la única cosa en Ceres que el Gremio controla mejor que las llaves.

Hubo alguna risa en la sala. La Directora no se rio, pero la comisura de su boca se movió un milímetro.

\escena
---Una última cosa ---dijo Loew, levantando un dedo---\nobreak. Querida, lo que usted ha hecho se parece muchísimo al contador partido de las colonias exteriores.

---Es el contador partido, profesor ---dijo Inés, sin dudarlo---\nobreak. Visto desde el otro lado. Ellos ponen en el cartel un contador que empieza en cero y sube con cada lector, y al final, cuando el cartel cambia, pasan ese contador a la hoja. Yo pongo un talonario que empieza lleno y baja con cada lector, y la hoja lo tiene apuntado desde el principio. En todo momento, mis tickets son exactamente $C$ menos su contador. Es la misma cuenta. La mía se paga por adelantado y la suya a plazos.

---Entonces, ¿qué ha inventado usted?

---Dos cosas pequeñas. Una forma de explicarlo en la que se ve de un vistazo que la hoja siempre va por arriba. Y la salida temprana: irse en cuanto se sabe que otro ya recargó. Los de las colonias solo se van si su propio cotejo funciona, o si ven el talonario alto. Por eso son la puerta giratoria. La salida temprana es lo que convierte la puerta giratoria en una puerta con turno.

---Dos cosas pequeñas ---repitió Loew, con una sonrisa---\nobreak. La historia de la ciencia está hecha de cosas pequeñas que nadie vio a tiempo.

La Directora Sørensen se puso de pie. Durante un momento, todos en la sala pensaron que iba a abrir la caja de llaves. En vez de eso, la cerró.

---La demostración me convence, doctora ---dijo---\nobreak. Pero yo no dirijo un aula. Dirijo un Archivo. Quiero ver su método funcionando con muchos lectores de verdad, contra la llave, contra el contador partido y contra los demás remedios que tenemos. En el Cerebro, a toda velocidad, millones de veces. Y si gana, el Gremio lo adoptará.

---¿Y si pierde?

---Si pierde, colgará usted misma la primera llave en el cartel del muelle doce. Lo prometió ayer.

---Lo prometí ---dijo Inés---\nobreak. Siete, avisa al Cerebro. Esta noche no vamos a dormir tampoco.
